\chapter{Граничный \texorpdfstring{$S_4\to C_4$}{S4 to C4}-компакттон и эта-диссипация}
\label{ch:v6-spectral-transition-discrete-compacton-c4-boundary-eta-dissipation-gate}

\section{Усиленная гипотеза}

Предыдущий гейт запретил канонический выбор $C_4$ внутри полного
аффинного $S_4$ и показал, что конечный когерентный сток не является
аттрактором. Здесь проверяется более сильная гипотеза, в которой два
прежних запрета снимаются явными физическими данными:
\begin{enumerate}
 \item внутреннее направление компакттона фиксируется граничным сектором;
 \item граничная голономия выбирает один четырёхцикл и ломает
 $S_4$ до его $C_4$-централизатора;
 \item слабая необратимость должна возникнуть из уже известной
 эта/Pfaffian-фазы, а не из нового произвольного резервуара.
\end{enumerate}

Антикруговая сверка использует
\path{version4_pfaffian_eta_orientation_gate},
\path{version4_determinant_line_inflow_gate},
\path{version5_eta_wzw_real_pair_phase_gate},
\path{version5_massless_holonomy_defect_index_gate} и
\path{version6_spectral_transition_compacton_c4_affine_selector_admissibility_gate}.

\section{Компакттон с выделенным внутренним направлением}

Выберем сбалансированный локальный вектор
\begin{equation}
 v_*=\left(\frac12,0,\frac12\right)\in\mathbb C^2\oplus\mathbb C.
 \label{eq:v6-c4-eta-selected-axis}
\end{equation}
Он удовлетворяет
\begin{equation}
 \|v_*\|^2=\frac12,
 \qquad
 \|v_{*,L}\|^2=|v_{*,R}|^2=\frac14.
\end{equation}
Поэтому при $\kappa=2\pi$ две двухузловые конфигурации
\begin{equation}
 \Psi_+=(v_*,-iv_*),
 \qquad
 \Psi_-=(v_*,+iv_*)
 \label{eq:v6-c4-eta-selected-compactons}
\end{equation}
являются точными собственными состояниями редуцированного шага с
собственными значениями $+i$ и $-i$. Машинные остатки обоих равенств
нулевые.

Это снимает внутреннюю неединственность условно: ось теперь задана
граничным сектором. Настоящий гейт не выводит $v_*$ из вакуумной
динамики и не объявляет этот выбор спонтанным.

\section{Граница действительно оставляет ровно один \texorpdfstring{$C_4$}{C4}}

Пусть на аффинном четырёхточечном носителе задано условие
\begin{equation}
 \psi(s+L)=C\psi(s),
 \qquad
 C=(1234).
 \label{eq:v6-c4-eta-boundary-condition}
\end{equation}
Остаточная симметрия состоит из элементов $g\in S_4$, сохраняющих
граничное условие, то есть удовлетворяющих $gC=Cg$.

Полный перебор двадцати четырёх перестановок даёт
\begin{equation}
 Z_{S_4}(C)=\{1,C,C^2,C^3\}=\langle C\rangle\simeq C_4.
 \label{eq:v6-c4-eta-centralizer}
\end{equation}

\begin{proposition}[Точное граничное нарушение]
Граничное условие~\eqref{eq:v6-c4-eta-boundary-condition} математически
ломает $S_4$ ровно до одной выбранной подгруппы $C_4$. На аффинном
триплете остаются три линии характеров $i,-1,-i$, причём линии $\pm i$
лежат в $P_3\mathbb C^4$ и могут быть сопоставлены двум компакттонам
\eqref{eq:v6-c4-eta-selected-compactons}.
\end{proposition}

Таким образом, первая половина предложенного механизма проходит. Цена
прохода ясна: конкретный четырёхцикл является граничным входом, а не
следствием полного $S_4$-родителя.

\section{Эта-фаза выбранных характеров}

Для одномерного оператора Дирака со скручиванием
$e^{2\pi i\alpha}$ спектр пропорционален $n+\alpha$, а при
$0<\alpha<1$ приведённый спектральный дисбаланс имеет форму
\begin{equation}
 \eta_\alpha(0)=1-2\alpha.
 \label{eq:v6-c4-eta-twisted-circle}
\end{equation}
Для характеров $+i$ и $-i$ получаем
\begin{equation}
 \alpha_+=\frac14,
 \quad
 \alpha_-=\frac34,
 \qquad
 \eta_+=\frac12,
 \quad
 \eta_-=-\frac12.
\end{equation}
Стандартная эта-фаза фермионного интеграла равна
\begin{equation}
 Z_{\eta,\pm}
 =\exp\!\left(-\frac{i\pi}{2}\eta_\pm\right)
 =e^{\mp i\pi/4}.
 \label{eq:v6-c4-eta-unit-phases}
\end{equation}
Она действительно различает две линии: относительная фаза равна $-i$.
Но
\begin{equation}
 |Z_{\eta,+}|=|Z_{\eta,-}|=1,
 \qquad
 -\log|Z_{\eta,\pm}|=0.
 \label{eq:v6-c4-eta-zero-decay}
\end{equation}

Тот же вывод следует для reduced Pfaffian-знаков $\pm1$: их модули
единичны. В полном KO6-блоке две Real-сопряжённые фазы дополнительно
сокращаются до $+1$, как уже доказано в Томах IV--V.

\begin{theorem}[Фаза не является скоростью распада]
Выбранная $C_4$-граница делает эта-инварианты двух компакттонных
характеров противоположными, но эта- и Pfaffian-вклады остаются чистыми
фазами. Они не имеют отрицательной вещественной части эффективного
действия и потому сами не задают положительную скорость диссипации.
\end{theorem}

Эта граница согласуется с общей геометрией фермионного интеграла, где
эта-инвариант задаёт фазу и ориентацию Pfaffian/determinant line
\cite{c4-eta-witten2015}, тогда как необратимая динамика требует отдельной
полностью положительной полугруппы
\cite{c4-eta-gks1976,c4-eta-lindblad1976}.

\section{Три минимальные редуцированные динамики}

В базисе $|+i\rangle,|-i\rangle$ положим
\begin{equation}
 Q_\eta=P_{+i}-P_{-i}=\operatorname{diag}(1,-1).
\end{equation}

\subsection{Когерентная эта-фаза}

Гамильтониан $H_\eta=Q_\eta/2$ вращает относительную фазу, но сохраняет
обе населённости $w_\pm$. Поэтому дефект характерной чистоты
$4w_+w_-$ постоянен и захвата нет.

\subsection{Минимальная эта-дефазировка}

Даже если феноменологически превратить тот же диагональный оператор в
скачок
\begin{equation}
 L_\varphi=\sqrt\gamma\,Q_\eta,
 \end{equation}
получается только дефазировка. При исходном
$w_+=w_-=1/2$ когерентность убывает к нулю, чистота убывает к $1/2$, но
\begin{equation}
 w_+(t)=w_-(t)=\frac12,
 \qquad
 4w_+(t)w_-(t)=1.
\end{equation}
То есть система приходит к равной классической смеси двух компакттонов,
а не выбирает один из них.

\subsection{Ориентированный амплитудный скачок}

Настоящий захват возникает для
\begin{equation}
 L_\downarrow=\sqrt\gamma\,|+i\rangle\langle-i|.
 \label{eq:v6-c4-eta-oriented-jump}
\end{equation}
Тогда
\begin{equation}
 w_+(t)=1-(1-w_+(0))e^{-\gamma t},
 \qquad
 w_-(t)=w_-(0)e^{-\gamma t},
 \label{eq:v6-c4-eta-capture-law}
\end{equation}
и характерный дефект стремится к нулю. При $\gamma=10^{-2}$ и
$w_+(0)=1/2$ аудит получает $w_+(500)=0.996631\ldots$.

Однако оператор~\eqref{eq:v6-c4-eta-oriented-jump} не является функцией
$Q_\eta$: алгебра эта-данных
$\operatorname{span}\{I,Q_\eta\}$ диагональна, тогда как
$L_\downarrow$ является внедиагональной частичной изометрией. Остаток его
проекции на эту алгебру равен $0.1$ при тестовой
$\gamma=10^{-2}$. Кроме того, время достижения $w_+=0.99$ равно
\begin{equation}
 t_{0.99}=\frac{\log50}{\gamma};
\end{equation}
оно произвольно меняется с невыведенной $\gamma$.

\begin{proposition}[Условный захват]
После выбора оси, ориентации границы и отдельного неэрмитова скачка
$L_\downarrow$ редуцированный компакттон экспоненциально захватывается в
ветвь $+i$. Но эта/Pfaffian-фаза поставляет только диагональную ориентацию
характеров; она не выводит ни частичную изометрию
$|-i\rangle\to|+i\rangle$, ни положительное число $\gamma$.
\end{proposition}

\section{Область действия результата}

Даже условный скачок действует только внутри уже выделенного точного
компакттонного многообразия. Он не возвращает пространственно вытекшую
амплитуду, не создаёт хиральный баланс и не захватывает обычный
одноузловой или гауссов пакет. Поэтому он решает фазовую неединственность
после подготовки компакттона, но не эндогенное рождение компакттона из
вакуума.

\begin{nogocriterion}[Запрет превращения фазы в трение]
Нельзя отождествлять мнимую эта/Pfaffian-фазу с малой положительной
скоростью. Для переоткрытия необходимо вывести из конкретного
граничного фермионного оператора спектральную плотность резервуара,
внедиагональную частичную изометрию между $C_4$-характерами и её
положительный коэффициент. Аналитическое продолжение фазы в затухание без
такого вывода считается новым феноменологическим правилом.
\end{nogocriterion}

\section{Вердикт}

Предложенная комбинация разделяется на точный положительный и строгий
отрицательный результаты. Выделенная ось и граничное условие
$S_4\to C_4$ дают непротиворечивый компакттонный сектор. Если отдельно
вставить слабый ориентированный скачок, он действительно создаёт
экспоненциальный фазовый захват.

Но эта- и Pfaffian-фазы имеют единичный модуль, полный KO6-пакет сокращает
их ориентацию, а доступная диагональная алгебра даёт максимум
дефазировку. Следовательно, слабая диссипация из этих фаз не выведена.

Следующий гейт
\path{version6_spectral_transition_discrete_compacton_energy_degeneracy_boundary_overlap_gate}
должен проверить, не возникает ли предпочтение без нового $\gamma$ из
разности вещественных энергий, нелинейности профиля или граничного
матричного элемента.

\begin{protocol}
\begin{enumerate}
 \item выделенная сбалансированная ось: \textbf{совместима};
 \item точные компакттоны $\pm i$ на этой оси: \textbf{получены};
 \item централизатор выбранного четырёхцикла: \textbf{ровно $C_4$};
 \item граничное нарушение $S_4\to C_4$: \textbf{точное};
 \item происхождение самой границы из родителя: \textbf{не выведено};
 \item эта-инварианты линий: \textbf{$+1/2,-1/2$};
 \item относительная эта-фаза: \textbf{$-i$};
 \item скорость из модуля эта-фазы: \textbf{ноль};
 \item полный KO6 Pfaffian: \textbf{фаза сокращается};
 \item диагональная эта-дефазировка выбирает ветвь: \textbf{нет};
 \item вставленный скачок $|+i\rangle\langle-i|$ захватывает:
 \textbf{да};
 \item оператор скачка из эта-алгебры: \textbf{нет};
 \item коэффициент $\gamma$: \textbf{свободен};
 \item захват обычных локализованных состояний: \textbf{не получен};
 \item R2, эндогенный trigger: \textbf{нет};
 \item R3, скорость: \textbf{нет};
 \item R4: \textbf{условно внутри заранее подготовленного сектора}.
\end{enumerate}
\end{protocol}

Машинный сертификат сохранён в
\path{s2t_v6_spectral_transition_discrete_compacton_c4_boundary_eta_dissipation_gate_results.json}.

\begin{thebibliography}{9}
\bibitem{c4-eta-witten2015}
E.~Witten,
\emph{Fermion Path Integrals and Topological Phases},
Reviews of Modern Physics 88 (2016), 035001,
\path{arXiv:1508.04715}.

\bibitem{c4-eta-gks1976}
V.~Gorini, A.~Kossakowski, and E.~C.~G.~Sudarshan,
\emph{Completely positive dynamical semigroups of N-level systems},
Journal of Mathematical Physics 17 (1976), 821--825,
\path{doi:10.1063/1.522979}.

\bibitem{c4-eta-lindblad1976}
G.~Lindblad,
\emph{On the generators of quantum dynamical semigroups},
Communications in Mathematical Physics 48 (1976), 119--130,
\path{doi:10.1007/BF01608499}.
\end{thebibliography}